% Source: https://artofproblemsolving.com/community/c3h1688279_hello_question
% Stable identity: aops-topic-1688279-q01
For positive integers, $2^a$ is a product of $a$ twos and $2^b$ is a product of $b$ twos, so $2^a\cdot2^b$ is a product of $a+b$ twos:
\[
2^a\cdot2^b=\underbrace{2\cdots2}_{a}\cdot\underbrace{2\cdots2}_{b}=2^{a+b}=2^c .
\]
Example: $2^3\cdot2^2=8\cdot4=32=2^5$, and $3+2=5$.

For other integers, the definitions $2^0=1$ and $2^{-n}=\dfrac{1}{2^n}$ are chosen exactly so that the rule keeps working. If $a=0$: $2^0\cdot2^b=2^b=2^{0+b}$. If $a=-n<0$ and $b\ge n$: $2^{-n}\cdot2^b=\dfrac{2^b}{2^n}=2^{b-n}$, because $n$ of the $b$ twos cancel. If $a=-n$ and $0\le b<n$: $\dfrac{2^b}{2^n}=\dfrac{1}{2^{n-b}}=2^{b-n}$. If both are negative, $2^{-n}\cdot2^{-m}=\dfrac{1}{2^{n+m}}=2^{-(n+m)}$. In every case $2^a\cdot2^b=2^{a+b}$, so $a+b=c$ gives $\boxed{2^a\cdot2^b=2^c}$.
